\section{记录位点}\label{chap3_2:results_recording_sites}

在本课题中，ACC的记录范围对应BA 24区（扣带回背侧），OFC对应BA 11区，DLPFC对应BA 46区，各脑区的前后（AP）坐标范围以耳棒零点（ear bar zero）为参考，标注于图~\ref{fig2:recording_sites}。共记录到679个ACC神经元（猕猴W：247；猕猴T：432；猕猴U无ACC记录）、1056个DLPFC神经元（猕猴W：127；猕猴U：491；猕猴T：438）以及903个OFC神经元（猕猴W：111；猕猴U：446；猕猴T：346）。

\begin{figure}[H]
    \centering
    \includegraphics[width=\textwidth]{Figure2_RecordingSites}
    \bicaption{\enspace 记录位点}{\enspace Recording sites}
    \label{fig2:recording_sites}
\end{figure}
\fignotetext{ACC（红色）、DLPFC（绿色）和OFC（蓝色）的记录位点分别叠加标注在三只猕猴的冠状面结构磁共振图像上：猕猴W（左）、猕猴U（中）、猕猴T（右）。前后（AP）坐标以耳棒零点为参考，标注于各截面左侧。本实验使用线性阵列电极记录，每个矩形标记代表一个记录位点，矩形沿电极轴方向的长度约为2毫米，略大于实际记录范围（约1.5毫米），以计入记录过程中可能发生的电极漂移与定位误差；重复记录的位点仅绘制一次。箭头标注主要脑沟的解剖名称。}{Recording sites for ACC (red), DLPFC (green), and OFC (blue) were overlaid on coronal structural MRI images for each monkey: Monkey W (left), Monkey U (middle), and Monkey T (right). Anterior–posterior (AP) coordinates relative to ear bar zero were indicated above each section. Linear array electrodes were used; each rectangle represented one recording site, with the length along the electrode axis set to approximately 2~mm — slightly larger than the actual recording span ($\sim$1.5~mm) — to account for electrode drift and localization uncertainty. Sites recorded on multiple sessions were shown only once. Arrows indicated major sulci with their anatomical names.}
